\reportheading{7. Design comparison and conclusion}

\subhead{7.1 Comparison of calculation conditions}

DIALux tested three sensitivity scenarios: 90\% output, lower ceiling/wall/floor reflectances (70/50/20\%), and both combined, while retaining the original layout and MF. Table 4 shows the workstation performance.

\tabletitle{Table 4. Workstation results by calculation condition.}

\begin{reporttable}{L{0.34}L{0.25}L{0.2}L{0.21}}
\rowcolor{tablegrey} \textbf{Case} & \textbf{Lowest desk mean} & \textbf{Lowest \(U_{\mathrm{o}}\)} & \textbf{Highest UGR} \\ \hline
Requirement & ≥500 lx & ≥0.60 & ≤19 \\
Baseline & 593 lx & 0.80 & 18.1 \\
90\% output & 534 lx & 0.80 & 17.7 \\
Lower reflectance & 566 lx & 0.79 & 18.6 \\
Combined & 509 lx & 0.79 & 18.2 \\
\end{reporttable}

\subhead{7.2 Design choice}

At 90\% output, the lowest desk mean is 534 lx, 6.8\% above the target, and the worst directional UGR falls from 18.1 to 17.7. In contrast, lower reflectance at full output reduces the minimum desk mean to 566 lx while the UGR rises to 18.6, leaving a 0.4 margin below the limit.

With both changes combined, the lowest desk mean is 509 lx and the maximum directional UGR is 18.2. This relatively narrow 9 lx illuminance margin supports keeping full output as the design setting and tuning the installed system after completion.

\subhead{7.3 Maintenance and conclusion}

The maintenance factor is calculated as LLMF × LSF × LMF × RSMF = 0.90 × 1.00 × 0.95 × 0.94 = 0.8037, which is rounded down to 0.80 [4]. Specifically, the product declares L90 at 50,000 h [2], while the remaining factors assume prompt replacement, annual cleaning, and retained room finishes.

Full output remains the design setting, as the 593 lx minimum maintained desk mean is 19\% above the 500 lx target and corresponds to about 741 lx initially under the adopted MF. Accordingly, local dimming can be commissioned after installation against the actual finishes, desks, and screens, though dimmed power should be verified separately.
